\section{Nietzsche, Aquinas, Dante}

Recently a young friend announced to me his conversion from Tradition to Nietzsche's ``philosophy''. In conversation, he was rabid, like someone possessed by a gnome. I recognized that ``true believer'' feeling of satisfaction, very similar to the self-satisfied feeling of accomplishment after I take my first dump in the morning.

\paragraph{Esoteric vs Exoteric}


His commitment, and understanding, of Tradition was apparently shallow, since he was unable to grasp some basic distinctions, much like the way some Spanish speakers I know have trouble hearing the difference between the sound of ``b'' and ``v''. The first distinction is between the exoteric and esoteric aspects of a religious Tradition. I am simply not interested in debating exoterisms or defending every opinion or event of the past 1500 years of Western Civilization. He would toss out poorly understood Bible quotes like some TV preacher, so I was a little surprised when he didn't ask for donations. Yet, like the preacher, he expected to acquire a convert.

\paragraph{Psychological vs Spiritual}


The second distinction is the confusion between the psychic and the spiritual, something Guenon warned against as one of the signs of the times. So my friend proved himself not to be opposed to the times, but rather the unconscious instrument of further decadence. Nietzsche's claim to fame is as a psychologist. In that, he can be very insightful in exposing hypocrisy, ressentiment, and the like. In that as a side effect, he is a creator of the modern world which has replaced discussion and dialog with accusation. Today, a man's ideas are barely considered, but instead his motives are questioned; certain points of view are even considered to be ``mental illnesses''. Nietzsche, and unfortunately my friend, have no awareness of spiritual states or higher states of being. This makes a fruitful discussion simply impossible.

\paragraph{Salvation vs Deliverance}


Another distinction is that between salvation and deliverance, as Guenon often points out. Salvation is offered to all, although not necessarily effectively. Deliverance is for the few. When reading religious texts, this distinction must be kept in mind, as it may not always be explicitly made.

Since my friend was unable to grasp these distinctions, he left unconvinced. Nevertheless, he sent me a passage from Nietzsche's \emph{Genealogy of Morals}, as though it was game, set, and match for his cause. I will briefly address the comments on Aquinas and Dante, since much of the conversation centered on them, saving a fuller explanation for another time.

\paragraph{Aquinas on Hell}


Nietzsche misrepresents Aquinas' alleged claim that the saints take pleasure in the suffering of the damned. An intellectually honest man will go to the source in the \emph{Summa Theologica\footnote{\url{https://www.newadvent.org/summa/5094.htm}}}. As I tried to explain to deaf ears, the rejoicing comes from the idea of justice. Justice the highest virtue so joy should accompany it. My friend couldn't object to justice per se, yet his mind could not grasp this at all. As an anti-egalitarian, as he claimed to be, he should know that some men will sink lower and others rise higher; unlike the modern mind, this is seen to be the nature of things, not an injustice. If you meditate on that passage, you may come to a direct insight that it is true; but that is something every man must do for himself, as my personal claim has no authority over you.

As for the notion of hell itself, this is hardly an innovation that arrived with Christianity. It is an idea common to Tradition including the Indian (Hinduism, Buddhism, Jainism), Islam, and even pagans; Valhalla was reserved for the heroes. But this is not the time for a fuller discussion, other than to point out that esoterically heaven and hell represent states of being. Far from being anti-life, understanding these states, and progressing \emph{spiritually}, not psychologically, \textbf{in this life}, is the path to fulfillment.


\paragraph{Understanding Dante}


Dante is a brilliant poet whose \emph{Divine Comedy} can be understood on several levels: as a work of art, as a socio-political commentary, as theology, or as an esoteric path. Dante is difficult to understand and 800 years of commentary has not exhausted the topic. Nietzsche, on the other hand, is easy to understand, and thus is popular with teenage boys. My friend believes that Nietzsche has done what 800 years of study had been unable to achieve: to dismiss Dante with an ill-conceived psychological diagnosis.

Does anyone reading Gornahoor seriously believe that Dante was writing a National Geographic style guidebook to Hell? Please raise your hand and reassure me if not, as I am feeling depressed about this conversation, at least insofar as I can still be bother with such human, all-too-human sentiments.

Quickly, Dante portrays the three stages of the spiritual path, which is not the same as becoming psychologically healthy minded.

\begin{enumerate}


\item The first stage is the natural state of the philosopher, usually called \textbf{ataraxia} or apatheia.

\item The next stage, having mounted Purgatory, is the Primordial State, corresponding to \textbf{salvation}.

\item Finally, there is the Supreme Identity, or \textbf{deliverance}.
\end{enumerate}


Has Nietzsche demonstrated any understanding of that? No, instead, he wants to accuse Dante of envy for relegating the noble pagans to Limbo. Now Dante was not attacking the pagans, as he shows high regard for them. Rather, he used them as symbols, as is obvious to anyone who knows how to read esoteric literature. His purpose is to show that the goal of ataraxia is far from what is possible for most men to achieve \emph{in this life}. Furthermore, this path is not exclusively Christian, as Dante's poem is an elaboration of similar Muslim ideas such as Ibn Arabi's, or Muhammed's own visit to heaven. The Indian religions also teach that a man must overcome certain spiritually destructive habits of mind, much like the path through Purgatory.

The idea of detachment from life is not anti-life as we pointed out several years ago in \textit{Natural and Supernatural Happiness}\footnote{\url{https://gornahoor.net/?p=13}}. A man can take this heroic path, or he can wait patiently for the eternal return of the same old.

\paragraph{Genealogy of Morals}


\begin{quotationx}
In belief in what? In love with what? In hope for what?---There's no doubt that these weak people---at some time or another they also want to be the strong people, some day their ``kingdom'' is to arrive---they call it simply ``the kingdom of God,'' as I mentioned. People are indeed so humble about everything! Only to experience that, one has to live a long time, beyond death---in fact, people must have an eternal life, so they can also win eternal recompense in the ``kingdom of God'' for that earthly life ``in faith, in love, in hope.'' Recompense for what? Recompense through what?

... In my view, Dante was grossly in error when, with an ingenuity inspiring terror, he set that inscription over the gateway into his hell: ``Eternal love also created me.'' Over the gateway into the Christian paradise and its ``eternal blessedness'' it would, in any event, be more fitting to let the inscription stand ``Eternal hate also created me''---provided it's all right to set a truth over the gateway to a lie! For what is the bliss of that paradise? ... Perhaps we might have guessed that already, but it is better for it to be expressly described for us by an authority we cannot underestimate in such matters, Thomas Aquinas, the great teacher and saint: ``In the kingdom of heaven'' he says as gently as a lamb, ``the blessed will see the punishment of the damned, so that they will derive all the more pleasure from their heavenly bliss.''

Or do you want to hear that message in a stronger tone, something from the mouth of a triumphant father of the church, who warns his Christians against the cruel sensuality of the public spectacles. But why? ``Faith, in fact, offers much more to us,'' he says (in de Spectaculis, c. 29 ff), ``something much stronger. Thanks to the redemption, very different joys are ours to command; in place of the athletes, we have our martyrs. If we want blood, well, we have the blood of Christ ... But what awaits us on the day of his coming again, his triumph!''---and now he takes off, the rapturous visionary:

However there are other spectacles---that last eternal day of judgment, ignored by nations, derided by them, when the accumulation of the years and all the many things which they produced will be burned in a single fire. What a broad spectacle then appears! How I will be lost in admiration! How I will laugh! How I will rejoice! I will be full of exaltation then as I see so many great kings who by public report were accepted into heaven groaning in the deepest darkness with Jove himself and alongside those very men who testified on their behalf! They will include governors of provinces who persecuted the name of our Lord burning in flames more fierce than those with which they proudly raged against the Christians! And those wise philosophers who earlier convinced their disciples that God was irrelevant and who claimed either that there is no such thing as a soul or that our souls would not return to their original bodies will be ashamed as they burn in the conflagration with those very disciples! And the poets will be there, shaking with fear, not in front of the tribunal of Rhadamanthus or Minos, but of the Christ they did not anticipate!

Then it will be easier to hear the tragic actors, because their voices will be more resonant in their own calamity'' (better voices since they will be screaming in greater terror). ``The actors will then be easier to recognize, for the fire will make them much more agile. Then the charioteer will be on show, all red in a wheel of fire, and the athletes will be visible, thrown, not in the gymnasium, but in the fire, unless I have no wish to look at their bodies then, so that I can more readily cast an insatiable gaze on those who raged against our Lord. `This is the man,' I will say, `the son of a workman or a prostitute'\,'' (in everything that follows and especially in the well-known description of the mother of Jesus from the Talmud, Tertullian from this point on is referring to the Jews) ``the destroyer of the Sabbath, the Samaritan possessed by the devil. He is the man whom you brought from Judas, the man who was beaten with a reed and with fists, reviled with spit, who was given gall and vinegar to drink. He is the man whom his disciples took away in secret, so that it could be said that he was resurrected, or whom the gardener took away, so that the crowd of visitors would not harm his lettuce.' What praetor or consul or quaestor or priest will from his own generosity grant this to you so that you may see such sights, so that you can exult in such things? And yet we already have these things to a certain extent through faith, represented to us by the imagining spirit. Besides, what sorts of things has the eye not seen or the ear not heard and what sorts of things have not arisen in the human heart?'' (1. Cor. 2, 9). ``I believe these are more pleasing than the race track and the circus and both enclosures'' (first and fourth tier of seats or, according to others, the comic and tragic stages). Through faith: that's how it's written.
\end{quotationx}



\flrightit{Posted on 2012-07-17 by Cologero}

\begin{center}* * *\end{center}

\begin{footnotesize}
\begin{sffamily}
\texttt{Aghorable on 2012-07-17 at 06:41 said:}

Be reassured, good Sir; long it is since our release from the deep dungeons of the literalist labyrinth. With bad behavior, we were pardoned early, and adjusted to a richer life on the surface.

And good health to you! May your morning dumps continue to be smooth and thorough. A wise man's bowels purge more than the remains of last night's dinner, and Zarathustra knew this (``Das Leben ist ein Born der Lust... '').

And yes, if any wonder, despite my name, I do flush the toilet after my daily evacuation!

\hfill

\texttt{golgonooza on 2012-07-17 at 06:47 said:}

*raises hand* I don't believe Dante was writing a guidebook to Hell! I can see why the conversation got depressing. It's hard to argue with literalists of any sort. Reading the part in the Summa about what the blessed will see is interesting and sort of shocking to those who believe in universal salvation. Dante is clear in the Divine Comedy though that the different states that people find themselves are direct reflections of the sin they have committed -- he is clear about the way in which the sin itself locks the person into a prison of their own making. I'm also interested in the difference between salvation and deliverance -- in what way would they be different? Are we talking about the difference between say, those who will be purified in Purgatory and thus saved from the torments of Hell, and those who attain to the beatific vision?

\hfill

\texttt{Cologero on 2012-07-17 at 07:15 said:}

Yes, Aghorable and golgonooza, and I think there is a lesson here about how to read Nietzsche, again not too literally. As a polemicist, trying to grab the attention of a world too busy to listen, he often had to go ``over the top''. It takes a careful reading to determine exactly who and what he is criticizing.


\hfill

\end{sffamily}
\end{footnotesize}

